\documentclass{beamer}
\usetheme{Boadilla}
\usecolortheme{dove}
\setbeamertemplate{navigation symbols}{}
\usepackage{graphicx}
\usepackage{amsmath,amssymb}
\usepackage[numbers]{natbib}
\bibliographystyle{apalike}

\title{The Impact of Health Litigation on Public Procurement:\\ Theory and Evidence from Brazil}
\author{Your Name}
\institute{Your Institution}
\date{\today}

\begin{document}

\begin{frame}
  \titlepage
\end{frame}

\AtBeginSection[]{
  \begin{frame}
    \frametitle{Outline}
    \tableofcontents[currentsection]
  \end{frame}
}

\section{Motivation}

\begin{frame}{Motivation}
\begin{itemize}
  \item Universal access to healthcare is a constitutional right in Brazil, but gaps in public provision remain.
  \item Patients often resort to lawsuits to obtain medicines or treatments not readily provided by the public health system.
  \item Such judicial enforcement of health rights ensures individual access, but may disrupt planning and efficiency of public procurement.
  \item Evidence suggests that litigation for health services has been rising and creating significant financial and administrative pressures on the health system \citep{Vieira2007, TCU2017}.
\end{itemize}
\end{frame}

\begin{frame}{Background Literature}
\begin{itemize}
  \item \textbf{Equity concerns:} Early studies found that judicialization can deepen inequalities in access. For example, lawsuits in S\~ao Paulo tended to benefit better-off groups, raising fairness issues \citep{Chieffi2009}.
  \item \textbf{Cost impact:} Court-mandated health purchases often come at higher prices, straining budgets \citep{Vieira2007, Chieffi2009}.
  \item \textbf{Systemic effects:} By prioritizing court orders, health authorities may divert resources from rational planning \citep{Vargas2019}. Non-compliance with orders can lead to fines or even jail for officials \citep{Delduque2013}.
  \item These findings motivate our analysis of how litigation influences procurement outcomes and behavior of public officials.
\end{itemize}
\end{frame}

\begin{frame}{Rising Health Litigation in Brazil}
\begin{itemize}
  \item The number of lawsuits demanding medicines or treatments has surged in recent years.
  \item Budget expenditures on court-mandated health purchases have also grown, indicating an increasing reliance on the judiciary to fulfill health needs.
\end{itemize}
\begin{center}
  \includegraphics[width=0.8\textwidth]{figures/figure1.png}
\end{center}
\end{frame}

\begin{frame}{Research Questions}
\begin{itemize}
  \item \textbf{Q1:} How does the possibility of litigation affect the decision-making of public officials in charge of procurement?
  \item \textbf{Q2:} What are the differences in outcomes (cost, timeliness, etc.) between ordinary planned procurement and court-mandated (litigated) procurement?
  \item \textbf{Q3:} Do public health authorities adjust their procurement strategy in anticipation of (or response to) litigation threats?
\end{itemize}
\end{frame}

\begin{frame}{This Study's Contribution}
\begin{itemize}
  \item We develop a simple microeconomic model of a public official’s procurement choice under the threat of legal enforcement (litigation).
  \item We use novel administrative data on health procurement and lawsuits to provide quantitative evidence on the cost and efficiency trade-offs imposed by litigation.
  \item We document the magnitude of inefficiencies (higher prices, small purchase sizes) resulting from judicial intervention, and examine whether officials adapt behavior to mitigate these effects.
  \item Policy implications are drawn to improve procurement processes and reduce the need for litigation.
\end{itemize}
\end{frame}

\section{Institutional Context}

\begin{frame}{Legal Framework}
\begin{itemize}
  \item \textbf{Right to Health:} The 1988 Brazilian Constitution (Art. 196) establishes health as a right of all and a duty of the State. Citizens can invoke this right in court to obtain treatments.
  \item \textbf{Public Procurement Law:} Procurements are governed by Law 8,666/93, which mandates competitive bidding for most purchases, but provides exceptions for certain situations (e.g. emergencies, small amounts).
  \item \textbf{Judicialization:} Courts generally uphold the individual's right to treatment. Judges can issue injunctions ordering the government to procure specific medicines, enforceable by law.
  \item These legal provisions set the stage for three procurement routes in practice: ordinary bidding, administrative exceptions, and court-mandated purchases.
\end{itemize}
\end{frame}

\begin{frame}{Ordinary Procurement Process}
\begin{itemize}
  \item \textbf{Planned Budget Purchases:} Most health supplies are acquired through standard competitive bidding using planned budgets.
  \item \textbf{Procedure:} Public agencies solicit bids or conduct electronic auctions for bulk quantities of medicines/equipment on a regular schedule.
  \item \textbf{Timeline:} The process typically takes considerable time (often 30--180 days from requisition to delivery) due to advertising, bidding, and contracting steps.
  \item \textbf{Advantages:} Bulk purchasing and competition usually result in lower unit prices. The process is transparent and follows procurement regulations.
  \item \textbf{Drawback:} Less flexibility to respond quickly to urgent or unanticipated needs (rigid schedule and bureaucracy).
\end{itemize}
\end{frame}

\begin{frame}{Administrative Procurement (Exceptions)}
\begin{itemize}
  \item \textbf{Direct Purchase Exceptions:} Procurement law allows \emph{dispensa de licita\c{c}\~ao} (waiver of bidding) in special cases (e.g. emergencies, sole suppliers).
  \item \textbf{Funding:} Often financed through extra-budgetary funds or emergency funds, rather than the normal planned budget.
  \item \textbf{Characteristics:} Smaller quantities, targeted to immediate needs. The agency can directly contract a supplier.
  \item \textbf{Timeline:} Much faster (on the order of days to a couple of weeks) since it bypasses formal bidding procedures.
  \item \textbf{Accountability:} No legal penalties for using this route if justified (it's legally permitted), but it must be properly documented to avoid audit issues.
\end{itemize}
\end{frame}

\begin{frame}{Court-Mandated (Litigated) Procurement}
\begin{itemize}
  \item \textbf{Trigger:} Initiated by a court order in response to a citizen’s lawsuit demanding a specific health service or product not (timely) provided by the system.
  \item \textbf{Funding:} Generally paid from extra-budgetary or contingency funds, since these needs were not budgeted or planned.
  \item \textbf{Characteristics:} Typically very small quantities (often an individual patient’s requirement). The purchase is tailor-made for the case.
  \item \textbf{Timeline:} Must be executed promptly (often within 1–10 days) as per the judge’s injunction, making speed paramount.
  \item \textbf{Enforcement:} Non-compliance can lead to contempt of court, fines, personal liability, or other penalties for officials \citep{Delduque2013}. Thus, there is a strong incentive to fulfill the order quickly.
\end{itemize}
\end{frame}

\begin{frame}{Comparison of Procurement Types}
\begin{center}
  \includegraphics[width=0.9\textwidth]{./Presentation/figures/figure3.pdf}
\end{center}
\begin{itemize}\footnotesize
  \item Ordinary procedures rely on planned budgets, larger quantities, and long lead times, with no direct punishment for delays.
  \item Administrative exceptions use extra-budget funds for quicker, small purchases when needed, without penalties for use.
  \item Litigated purchases also use extra-budget funds and are extremely fast, but introduce a threat of legal punishment if orders are not met.
\end{itemize}
\end{frame}

\section{Theoretical Model}

\begin{frame}{Model Setup}
\begin{itemize}
  \item Consider a public health manager responsible for procuring a certain treatment for a population of patients.
  \item Total potential demand is $N$ units (patients in need). The manager can decide to procure $q$ units through the ordinary process in advance.
  \item Timeline:
        \begin{enumerate}
            \item \textbf{Planning stage:} Manager chooses $q$ to acquire via regular procurement (at unit cost $c$).
            \item \textbf{Demand realization:} Actual need materializes. If demand exceeds $q$, the shortfall $N-q$ represents patients not initially covered.
            \item \textbf{Enforcement stage:} Unserved patients may seek judicial relief. Courts, if petitioned, order the manager to supply remaining units immediately.
        \end{enumerate}
  \item The manager aims to fulfill health needs while managing costs and avoiding penalties.
\end{itemize}
\end{frame}

\begin{frame}{Citizen’s Decision to Litigate}
\begin{itemize}
  \item Each unmet patient (beyond $q$) decides whether to file a lawsuit for the treatment.
  \item Let $\alpha$ represent the probability that a given unmet need results in a lawsuit. (This probability reflects factors like patient awareness, access to legal aid, and urgency of need.)
  \item If a patient sues and wins:
    \begin{itemize}
      \item The treatment will be provided via a court order (litigated procurement).
      \item The patient obtains the needed medicine (benefit).
      \item The manager must comply, incurring procurement costs and potential legal burden.
    \end{itemize}
  \item If the patient does not sue (probability $1-\alpha$), the need remains unfulfilled and the manager faces no immediate additional cost or penalty for that case.
  \item $\alpha$ thus captures the enforcement intensity of the right to health (higher $\alpha$ means unmet needs are more likely to translate into lawsuits).
\end{itemize}
\end{frame}

\begin{frame}{Government’s Procurement Problem}
\begin{itemize}
  \item The manager’s objective is to minimize total expected cost (for simplicity, assume a risk-neutral cost-minimizer).
  \item Costs come from:
    \begin{itemize}
      \item \textbf{Ordinary provision:} planning and purchasing $q$ units at cost $c$ each.
      \item \textbf{Litigation-driven provision:} for each of the $N - q$ unmet units, if a lawsuit occurs (with probability $\alpha$), the manager must procure that unit via expedited purchase at a higher unit cost $c_{\ell}$.
      \item \textbf{Penalty:} Additionally, each lawsuit may impose an extra penalty or legal cost $P$ (e.g., court fees, fines, or administrative overhead).
    \end{itemize}
  \item Thus, expected total cost as a function of $q$:
\end{itemize}
\[
    C(q) = c \cdot q \;+\; \alpha \cdot (c_{\ell} + P) \cdot (N - q).
\]
\begin{itemize}
  \item Here $c_{\ell} (> c)$ reflects the higher unit cost of ad-hoc litigated purchases.
  \item $P$ captures any additional per-case penalty or inefficiency from the litigation process.
\end{itemize}
\end{frame}

\begin{frame}{Optimal Solution}
To find the cost-minimizing $q^*$, take the derivative of $C(q)$ and set to zero (first-order condition):
\begin{align*}
    \frac{dC}{dq} &= c \;-\; \alpha (c_{\ell} + P) \;=\; 0. \\
    \implies c &= \alpha \,\big(c_{\ell} + P\big).
\end{align*}
\begin{itemize}
  \item The manager balances the marginal cost of procuring one more unit in advance ($c$) with the expected marginal cost of not procuring it and potentially facing a lawsuit [$\alpha(c_{\ell}+P)$].
  \item Solve for the interior solution:
    \[
       q^* = N \quad \text{if } c < \alpha(c_{\ell}+P), 
       \qquad 
       q^* = 0 \quad \text{if } c > \alpha(c_{\ell}+P).
    \]
  \item In other words, if the expected cost of a lawsuit for an unmet unit exceeds the cost of planning provision, the manager will fully supply the need (cover all $N$). If the lawsuit cost is very low relative to $c$, the manager might not stock any units (extreme case).
  \item For intermediate cases, this model yields a corner solution because of linear cost structure – in practice, the manager might choose partial coverage if marginal conditions just balance.
\end{itemize}
\end{frame}

\begin{frame}{Comparative Statics}
\begin{itemize}
  \item \textbf{Effect of litigation likelihood ($\alpha$):} An increase in the probability of lawsuits makes it more costly to under-provide. Thus, a higher $\alpha$ encourages the manager to plan for more units ($q^*$ increases).
  \item \textbf{Effect of penalty ($P$):} Stronger penalties or higher legal costs (fines, etc.) also raise $\alpha(c_{\ell}+P)$, leading to a larger $q^*$. When officials face serious consequences for non-provision, they are more inclined to procure in advance to avoid litigation.
  \item \textbf{Effect of cost differences ($c_{\ell}$ vs $c$):} A greater gap between expedited purchase cost $c_{\ell}$ and regular cost $c$ means litigation is expensive, pushing toward higher $q^*$. Conversely, if urgent purchases were not much costlier, the incentive to pre-purchase diminishes.
  \item \textbf{Budget constraint (not explicitly in model):} In reality, limited budgets may prevent covering all needs even if $\alpha(c_{\ell}+P) > c$. This can force $q < N$ despite high litigation costs, leading to some level of litigation as a relief valve.
\end{itemize}
\end{frame}

\begin{frame}{Implications of the Model}
\begin{itemize}
  \item The model suggests that litigation serves as a enforcement mechanism: if the threat is strong (high $\alpha$, big $P$), a rational manager will try to comply (supply more) to avoid being taken to court.
  \item However, if litigation is relatively infrequent or penalties are mild, the manager might accept some lawsuits as a cost trade-off, resulting in under-provision for those who do not sue.
  \item In equilibrium, some unmet demand and consequent lawsuits may occur if $c$ is high and $\alpha$ or $P$ are not high enough to warrant full coverage. This leads to a mix of ordinary and litigated procurement.
  \item The model highlights a potential inefficiency: the state may end up paying more per unit (at cost $c_{\ell}$) for those units acquired via litigation, which could have been cheaper at $c$ if planned.
  \item These predictions can be tested empirically: we expect litigated purchases to show higher unit costs and occur primarily when the public system fails to fully provide certain treatments.
\end{itemize}
\end{frame}

\section{Data}

\begin{frame}{Data and Sources}
\begin{itemize}
  \item \textbf{Procurement Records:} We compile administrative data on public health procurement transactions, including ordinary bids, administrative purchases, and court-mandated purchases. 
  \begin{itemize}
    \item Source: Health Secretariat procurement databases (years 20XX--20YY) and financial reports.
    \item Variables: item description, quantity, price paid, procurement modality (bidding, direct purchase, or judicial order), purchase date, etc.
  \end{itemize}
  \item \textbf{Lawsuit Records:} Data on health-related lawsuits obtained from the State Justice Department and Attorney General’s Office.
  \begin{itemize}
    \item We have records of court orders for medical treatments, including the date of court decision, the item requested, and compliance information.
  \end{itemize}
  \item We merge these sources to identify which procurement transactions were triggered by litigation and to measure outcomes for each type of purchase.
\end{itemize}
\end{frame}

\begin{frame}{Sample Overview}
\begin{itemize}
  \item Dataset covers **N =** XX,XXX procurement transactions over the study period.
  \item Of these, a certain fraction were court-mandated purchases:
    \begin{itemize}
      \item \textbf{Litigated purchases:} $Z$ transactions (about Y\% of all transactions) were explicitly linked to court orders.
      \item \textbf{Ordinary purchases:} The majority (around (100--Y)\%) were via planned tenders or routine processes.
      \item \textbf{Administrative (direct) purchases:} A smaller subset, roughly X\%, were emergency or exceptional purchases without bidding (non-judicial).
    \end{itemize}
  \item Total spending in the sample: approximately R\$ A million.
    \begin{itemize}
      \item Spending on court-mandated purchases accounts for roughly R\$ B million (B\% of total health procurement spending in the period).
      \item This indicates a non-trivial budget impact attributable to judicialization.
    \end{itemize}
  \item We observe the key outcomes of interest for each transaction:
    \begin{itemize}
      \item Unit price paid, quantity, and procurement lead time (time from need identification to delivery).
    \end{itemize}
\end{itemize}
\end{frame}

\begin{frame}{What Items Are Litigated?}
\begin{itemize}
  \item The majority of litigated court orders involve high-cost medications and treatments that are not readily available through the public system’s regular formulary.
  \item Common categories among litigated items:
    \begin{itemize}
      \item Specialized drugs for rare diseases (or orphan drugs) that are not provided by default due to cost or regulatory approval status.
      \item Oncology (cancer) medications and advanced therapies that patients demand before they are included in public protocols.
      \item High-cost medical supplies or nutritional supplements for specific chronic conditions.
    \end{itemize}
  \item By contrast, routine medications (e.g., basic antibiotics, hypertension drugs) are usually procured via ordinary channels and rarely appear in lawsuits.
  \item \textbf{Example:} Expensive biologic drugs for autoimmune diseases have been frequent subjects of litigation, as they often were not budgeted but patients won court orders to obtain them.
\end{itemize}
\end{frame}

\begin{frame}{Geographical Distribution of Cases}
\begin{itemize}
  \item There is significant variation in the incidence of health litigation across regions:
    \begin{itemize}
      \item More economically developed areas (e.g., the state capital and urban centers) tend to have higher numbers of lawsuits, possibly reflecting greater awareness and access to legal assistance.
      \item Some interior or rural regions show fewer cases, which could indicate unmet needs that are not translated into legal action, raising equity concerns.
    \end{itemize}
  \item The map/figure illustrates the number of health-related lawsuits (or court-mandated purchases) by municipality or region.
\end{itemize}
\begin{center}
  \includegraphics[width=0.8\textwidth]{./Presentation/figures/figure2.png}
\end{center}
\end{frame}

\section{Empirical Strategy}

\begin{frame}{Empirical Approach}
\begin{itemize}
  \item To quantify the impact of litigation on procurement outcomes, we compare similar purchases made via different routes:
    \begin{itemize}
      \item Litigated (court-ordered) purchases 
      \item Ordinary planned purchases (biddings)
    \end{itemize}
  \item Key outcome metrics:
    \begin{itemize}
      \item \textbf{Unit price} paid for the item.
      \item \textbf{Procurement speed} (time from need to delivery).
      \item \textbf{Quantity per transaction} (scale of purchase).
    \end{itemize}
  \item Identification challenge: Items subject to litigation might differ systematically (e.g., tend to be expensive or rare items). A naive comparison of all litigated vs non-litigated purchases could capture these inherent differences rather than the effect of the procurement mode.
  \item Strategy:
    \begin{itemize}
      \item Include item fixed effects to compare costs for the \emph{same product} procured via different routes.
      \item Control for other factors (year, supplier, etc.) that might affect price or speed.
      \item Essentially, we ask: if the same drug is obtained normally versus via a lawsuit, how do price and delivery time differ?
    \end{itemize}
\end{itemize}
\end{frame}

\begin{frame}{Regression Specification}
We estimate panel regressions exploiting variation across items and over time:
\begin{equation*}
    \text{Cost}_{ij} = \beta \cdot \text{Litigated}_{ij} + \gamma X_{ij} + \mu_{i} + \tau_{t} + \varepsilon_{ij},
\end{equation*}
where:
\begin{itemize}\itemsep5pt
  \item $\text{Cost}_{ij}$ is the unit price paid for item $i$ in purchase $j$.
  \item $\text{Litigated}_{ij}$ is an indicator that the purchase was via court order. $\beta$ captures the average price premium (or difference) for litigated purchases.
  \item $X_{ij}$ includes controls such as quantity, procurement year, and possibly supplier characteristics.
  \item $\mu_{i}$ are item (product) fixed effects, ensuring comparison within the same item.
  \item $\tau_{t}$ are time fixed effects (e.g., year or month) to account for inflation or temporal trends in prices.
\end{itemize}
We run analogous specifications for other outcomes:
\begin{itemize}
  \item Delivery time (days) as the dependent variable.
  \item Whether $\beta$ differs across subgroups (e.g., by drug type) to check heterogeneity.
\end{itemize}
\end{frame}

\section{Main Results}

\begin{frame}{Result 1: Unit Prices are Higher under Litigation}
\begin{itemize}
  \item We find a substantial price difference: purchases made under court order are significantly more expensive per unit than those made via ordinary procurement.
  \item On average, a given medicine obtained through litigation costs about **50\% more** than the same medicine purchased through competitive bidding.
  \item This difference persists even after accounting for product and time fixed effects, indicating it’s not just different items or inflation driving the gap.
  \item Possible reasons:
    \begin{itemize}
      \item Lack of competition/negotiation: The government often must accept the price from a single supplier on short notice.
      \item Loss of bulk discounts: Litigated orders are typically one-off small quantities.
      \item Urgency premium: Suppliers may charge more knowing the government is under a court mandate to buy immediately.
    \end{itemize}
\end{itemize}
\begin{center}
  \includegraphics[width=0.75\textwidth]{figures/figure4.png}
\end{center}
\end{frame}

\begin{frame}{Result 2: Faster Delivery through Litigation}
\begin{itemize}
  \item Court-mandated procurement drastically reduces delivery times:
    \begin{itemize}
      \item Ordinary planned purchases took a median of **X** days from request to delivery.
      \item Litigated purchases were fulfilled in a median of **Y** days (often within one or two weeks).
    \end{itemize}
  \item This reflects the priority given to judicial orders – bureaucratic steps are skipped or expedited to comply with the court.
  \item The time savings can be critical for patient health outcomes, especially for life-saving medications needed immediately.
  \item However, the speed comes at the cost of the higher prices noted. Essentially, the public sector “pays a premium” for urgency.
  \item These findings highlight a trade-off: the litigation route achieves timeliness that the normal process often cannot, albeit inefficiently.
\end{itemize}
\begin{center}
  \includegraphics[width=0.75\textwidth]{figures/figure5.png}
\end{center}
\end{frame}

\begin{frame}{Why Are Litigated Purchases Costly?}
\begin{itemize}
  \item \textbf{No Economies of Scale:} Court orders are typically for one patient or small quantities, so the government loses bulk purchasing power. In contrast, normal tenders aggregate demand (dozens or hundreds of units) to get volume discounts.
  \item \textbf{Emergency Procurement Conditions:} The urgent nature means less time to seek competitive offers. Often the agency must use a sole-source or buy at retail market rates to meet the deadline.
  \item \textbf{Supplier Leverage:} Suppliers know the government is compelled to procure the item quickly (by legal mandate). This urgency can reduce the government's bargaining position, potentially leading to higher quotes.
  \item \textbf{Administrative Overheads:} Rushed procurement might incur extra administrative costs (special handling, legal processing), effectively increasing the cost attributed to those units.
  \item These factors align with our model's parameters: $c_{\ell}$ (the effective cost of last-minute purchase) is much higher than $c$ (planned cost), explaining why ex-post these litigated units are expensive.
\end{itemize}
\end{frame}

\begin{frame}{Result 3: Little Evidence of Proactive Adaptation}
\begin{itemize}
  \item We examine whether health authorities adjust future procurement in response to past litigation:
    \begin{itemize}
      \item Do items that frequently triggered lawsuits eventually get added to regular procurement plans in larger quantities?
      \item Does the experience of costly litigation push the manager to increase $q$ (planned supply) over time for those items?
    \end{itemize}
  \item \textbf{Finding:} No strong evidence of a systematic proactive response.
    \begin{itemize}
      \item Many medicines continued to be acquired via lawsuits year after year, without significant increases in pre-planned provision.
      \item The share of budget spent on litigated purchases remained high or grew slightly over the sample period, rather than shrinking.
    \end{itemize}
  \item This suggests institutional or budgetary constraints impede adjustment. Despite the inefficiencies, agencies may be unable or slow to incorporate certain high-cost treatments into their standard procurement programs.
  \item In a few notable cases, after repeated litigation the government did incorporate an item (leading to fewer lawsuits subsequently), but these appear to be exceptions rather than the rule.
\end{itemize}
\end{frame}

\section{Robustness Tests}

\begin{frame}{Robustness Checks}
\begin{itemize}
  \item We perform several robustness checks to ensure our results are not driven by specific assumptions or subsets:
    \begin{enumerate}
      \item \textbf{Product fixed effects vs. no fixed effects:} Even without item fixed effects (just comparing overall averages), litigated purchases show a large price premium. With fixed effects, the premium remains, confirming it’s not due to different mix of products.
      \item \textbf{Subsample (Apples-to-Apples):} Restricting the analysis to only those products which were procured both via ordinary and via litigation at different times. Within this subset, the price and time differences persist, reinforcing our conclusions.
      \item \textbf{Controlling for quantity:} Adding quantity as a control (since litigated purchases are smaller) does not eliminate the price premium. This indicates that price differences are not solely due to quantity differences, but also the process.
      \item \textbf{Outliers:} Removing extreme cases (very expensive orphan drugs, etc.) yields similar patterns, so results are not driven by a few outlier items.
    \end{enumerate}
\end{itemize}
\end{frame}

\begin{frame}{Robustness: Controlling for Product Characteristics}
\begin{itemize}
  \item The inclusion of product fixed effects is a stringent test: it compares prices for the same item under different procurement modes.
  \item Our regression with item fixed effects shows a $\beta$ (Litigation coefficient) that remains positive and significant.
  \item In fact, for many individual products, one can compare the unit price paid in a normal bulk purchase vs. in a litigated purchase:
    \begin{itemize}
      \item e.g., A cancer drug: regular tender price \$100 per dose; litigated purchase price \$160 per dose (60\% higher).
      \item e.g., An enzyme replacement therapy: tender price \$5,000; litigated price \$7,000.
    \end{itemize}
  \item The figure below illustrates estimated price differences (with confidence intervals) for several high-volume items. All lie above zero, indicating consistently higher costs when procured via lawsuits.
\end{itemize}
\begin{center}
  \includegraphics[width=0.75\textwidth]{figures/figure6.png}
\end{center}
\end{frame}

\begin{frame}{Robustness: Alternative Specifications}
\begin{itemize}
  \item We also tried different modeling choices:
    \begin{itemize}
      \item Using \textbf{log(price)} as the outcome to interpret differences in percentage terms. Results were similar: litigated purchases have about $\approx 40$--$60\%$ higher prices on average.
      \item A \textbf{median regression} (quantile regression) to reduce influence of extreme values. The median effect of litigation on price was still significantly positive.
      \item Considering \textbf{total cost per case}: since litigated purchases are smaller quantity but higher unit price, we checked total cost per patient. We find that often the total cost is still higher due to price, despite fewer units.
      \item Looking at \textbf{different subsets}: for example, only medications vs. only medical equipment. Both subsets showed the same qualitative pattern of higher cost and faster delivery under litigation.
    \end{itemize}
  \item None of these alternative approaches overturn the main conclusions, giving us confidence in the robustness of our findings.
\end{itemize}
\begin{center}
  \includegraphics[width=0.75\textwidth]{figures/figure7.png}
\end{center}
\end{frame}

\section{Policy Implications}

\begin{frame}{Policy Implications (I) – Strengthen Planning and Provision}
\begin{itemize}
  \item \textbf{Improve Early Provision:} The government should proactively include highly demanded but currently uncovered treatments into the regular supply system. If certain expensive drugs are repeatedly requested via courts, incorporating them into public programs (or subsidy schemes) can reduce future litigation and possibly obtain better prices through bulk negotiation.
  \item \textbf{Increase Health Budgets for Key Areas:} Often the root cause of judicialization is underfunding of specific treatments. Allocating more resources to cover these needs (especially for rare diseases or new therapies) would preempt many lawsuits.
  \item \textbf{Flexibility in Procurement Rules:} Introduce more flexible procurement mechanisms for urgent or specialized needs, so that managers can respond quickly \emph{without} needing a lawsuit. For instance, frameworks contracts or faster bidding for certain categories could achieve much of the speed of litigation at lower cost.
  \item Encouragingly, when the government has acted on these lines, it showed results: e.g., after the health system started providing a previously litigated cancer drug, related lawsuits dropped by about 25\% the next year \citep{Borges2018}.
\end{itemize}
\end{frame}

\begin{frame}{Policy Implications (II) – Judicial System Reforms}
\begin{itemize}
  \item \textbf{Technical Support for Courts:} Establish dedicated technical chambers or expert committees (as some states have done) to advise judges on medical claims. This can filter out requests for less cost-effective treatments and promote alternatives, reducing unnecessary expenditure.
  \item \textbf{Mediation Mechanisms:} Prior to litigation, implement an administrative mediation process where patients’ requests are reviewed by health authorities quickly. Many cases could be resolved administratively (providing the drug if appropriate) without going to court.
  \item \textbf{Guidelines for Judges:} The National Council of Justice (CNJ) could issue stronger guidelines to standardize decisions, considering budget impact and available alternatives. If judges coordinate with health policy (e.g., require consultation of treatment protocols), it can reduce irrational outlier orders.
  \item These measures aim to align the judicial enforcement with public health planning, preserving the right to health while minimizing distortions.
\end{itemize}
\end{frame}

\begin{frame}{Policy Implications (III) – Long-Term Structural Solutions}
\begin{itemize}
  \item \textbf{Health Technology Assessment (HTA):} Accelerate the process of evaluating and incorporating new therapies into the public system. If effective new drugs are added more quickly, patients won’t need to sue for them. Strengthening bodies like CONITEC (which recommends new drug inclusion) is key.
  \item \textbf{Budgetary and Legal Reforms:} Consider earmarking funds for high-cost treatments or creating special state programs for rare diseases, so that these are handled within the system rather than through courts.
  \item \textbf{Public Awareness and Equity:} Ensure equitable access by publicizing administrative pathways. Often, only savvy, better-off patients use litigation. By creating help desks or outreach, all patients in need could access life-saving drugs either through a streamlined admin request or be informed about options, reducing the current bias.
  \item Ultimately, the goal is to reduce the reliance on litigation by meeting legitimate health needs through the health policy itself – making the extraordinary remedy of courts less necessary.
\end{itemize}
\end{frame}

\section{Conclusion}

\begin{frame}{Conclusion}
\begin{itemize}
  \item \textbf{Summary of Findings:} Judicial enforcement of the right to health has a dual effect. It succeeds in delivering medicines to patients quickly, but at a substantially higher cost per unit and in a manner that circumvents normal procurement efficiencies.
  \item \textbf{Efficiency Trade-off:} Our analysis quantified a significant “price of urgency” – the public sector pays a premium (on the order of 50\% higher prices) to fulfill court orders, reflecting lost economies of scale and emergency conditions.
  \item \textbf{Behavior of Officials:} In theory, a high threat of litigation should incentivize pre-emptive provision. In practice, we observed limited adjustment. Bureaucratic and budget constraints mean agencies often continue business-as-usual, resulting in recurring lawsuits.
  \item \textbf{Implications\:} While judicialization ensures individual rights are honored, it can undermine collective rationality in resource allocation\#:contentReference[oaicite:2]{index=2}. There is a need to integrate this feedback loop so that the health system learns and adapts to reduce inefficiencies.
  \item \textbf{Balancing Rights and Efficiency:} Policymakers must strive to uphold the right to health through better health policy and funding, rather than via court orders, to achieve more equitable and efficient outcomes.
\end{itemize}
\end{frame}

\begin{frame}{Limitations}
\begin{itemize}
  \item \textbf{Scope of Data:} Our data is from one state (or a subset of years). The situation may differ in other regions or at the federal level. Results should be generalized with caution.
  \item \textbf{Unobserved Health Outcomes:} We focus on costs and process efficiency. We do not directly measure patient health outcomes. It’s possible that faster access via litigation improved health in ways that justify some costs – a welfare analysis is outside our scope.
  \item \textbf{Endogeneity:} There could be unobserved factors behind which items get litigated (e.g., severity of illness or drug novelty) that correlate with costs. We addressed this with fixed effects and robustness checks, but residual bias might remain.
  \item \textbf{Dynamic Effects:} Our finding of little adaptation might be time-limited. Policy changes in recent years (e.g., new laws in 2024) could alter behaviors after our sample period. Continuous monitoring is needed.
  \item \textbf{Legal Considerations:} We have not delved into the legal merits of cases. Some litigation might be deemed unnecessary or abusive, while some is essential – treating them uniformly in data might oversimplify the issue.
\end{itemize}
\end{frame}

\begin{frame}{Future Research}
\begin{itemize}
  \item \textbf{Comparative Studies:} Extending analysis to other states or countries (e.g., Colombia’s experience with health litigation) to see if similar cost-speed trade-offs exist, or if different legal systems mitigate the inefficiencies.
  \item \textbf{Patient Outcomes:} Linking procurement data with patient health outcomes would help assess the ultimate impact of litigation – do patients who win lawsuits have better health/improved survival? And at what cost-effectiveness?
  \item \textbf{In-depth Case Studies:} Qualitative research on how health secretariats handle court orders internally could reveal opportunities for streamlining. Interviews could uncover why certain adaptations (like stocking often-litigated drugs) are not happening.
  \item \textbf{Policy Interventions:} Evaluate reforms such as technical committees for judges or pre-court mediation. If some states implemented these, one could use a difference-in-differences approach to quantify how they affect the number of lawsuits and procurement costs.
  \item Overall, further research can inform how to protect individual rights to health while maintaining sustainable and efficient public healthcare systems.
\end{itemize}
\end{frame}

\begin{frame}[allowframebreaks]{References}
\nocite{*}
\bibliography{References}
\end{frame}

\end{document}